% Generated draft: requires mathematical and source review.
\documentclass[11pt]{article}
\usepackage[margin=25mm]{geometry}
\usepackage{fontspec}
\setmainfont{Latin Modern Roman}
\usepackage{amsmath,amssymb,mathtools,unicode-math}
\setmathfont{Latin Modern Math}
\usepackage{enumitem,needspace,xurl,hyperref}
\setlength{\parindent}{0pt}
\setlength{\parskip}{5pt}
\setlength{\emergencystretch}{3em}
\newcommand{\N}{\mathbb{N}}
\newcommand{\Z}{\mathbb{Z}}
\newcommand{\Q}{\mathbb{Q}}
\newcommand{\R}{\mathbb{R}}
\newcommand{\C}{\mathbb{C}}
\newcommand{\F}{\mathbb{F}}
\newcommand{\A}{\mathbb{A}}
\newcommand{\PP}{\mathbb{P}}
\newcommand{\E}{\mathbb{E}}
\newcommand{\eps}{\varepsilon}
\newcommand{\dd}{\,\mathrm{d}}
\begin{document}
\section*{UnsolvedMath ID: 30001563}
\textbf{Source problem code:} catalogue entry
\subsection{Definitions and mathematical statement}
Let \(A=\{0,1,3,4\}\subset\Z\), with each letter interpreted as its integer value. Let \(A^*\) denote the finite words over \(A\), including the empty word, and define the morphism \(\sigma:A^*\to A^*\), extended by concatenation, by
\[
\sigma(0)=03,\qquad \sigma(1)=43,\qquad \sigma(3)=1,\qquad \sigma(4)=01.
\]
Let
\[
w=\lim_{k\to\infty}\sigma^k(0)=w_0w_1w_2\cdots
\]
be the one-sided infinite fixed point beginning with \(0\); the limit is understood through its nested finite prefixes. For a finite word \(x=a_1\cdots a_m\), define its length and letter sum by
\[
|x|=m,\qquad S(x)=\sum_{j=1}^{m}a_j.
\]
A factor is a contiguous finite block of \(w\).

\textbf{Conjecture (supplied attribution: Conjecture 48 of the cited workshop report).} There are no three nonempty words \(x_1,x_2,x_3\in A^*\) such that \(x_1x_2x_3\) is a factor of \(w\) and
\[
|x_1|=|x_2|=|x_3|,\qquad S(x_1)=S(x_2)=S(x_3).
\]
Equivalently, for every integer \(i\geq0\) and every integer \(n\geq1\), the three integers
\[
\sum_{j=0}^{n-1}w_{i+j},\qquad
\sum_{j=0}^{n-1}w_{i+n+j},\qquad
\sum_{j=0}^{n-1}w_{i+2n+j}
\]
are not all equal.

The nonempty-block condition uses the usual convention for cube avoidance; the supplied formulas do not explicitly state it. Only equal lengths and equal integer letter sums are required: the blocks need not be identical or have identical letter multiplicities. The source labels this claim open and its transcription corrected and verified; those are supplied curation labels, not independent verification here.
\subsection{Short English statement}
Starting with \(0\), repeatedly apply the given substitution to obtain an infinite word. The conjecture says that no contiguous block splits into three nonempty consecutive pieces of equal length whose integer letter sums are all equal. The supplied open-status and source-verification assessments require independent review.
\subsection{Sources}
\begin{itemize}
\item Local registry: \texttt{data/unsolvedmath/problems.json}, supplied record with permanent UnsolvedMath ID \(30001563\), source problem code \texttt{catalogue entry}, titled ``PVHH-Cube Avoidance in a Morphic Fixed Point.'' The supplied record belongs to the Oberwolfach Reports Open Problems collection. No individual proposer is supplied.
\item Supplied source citation: \emph{Mini-Workshop: Combinatorics on Words (2011)}, Oberwolfach Reports, pp.\ 2195--2244, DOI \texttt{10.4171/owr/2010/37}. Original workshop report: \url{https://doi.org/10.4171/owr/2010/37}. The supplied statement-verification field attributes the morphism and equal-length, equal-sum cube-avoidance claim to Conjecture 48 and identifies this record as the canonical record for a duplicated extraction. This attribution has not been independently checked here. The record gives proposed year 2010, whereas its report citation says 2011; this discrepancy remains unresolved.
\item Additional supplied literature source, labeled ``Au, Robertson, and Shallit'': \url{https://arxiv.org/abs/0812.2466}. The record supplies no title or specific passage establishing its relationship to this conjecture.
\item Supplied literature assessment, checked on \texttt{2026-08-21}: the record labels the claim open and states that the extracted conjectural PVHH-cube claim remains unresolved, with available later literature providing no proof or counterexample for the stated fragment. The listed sources checked are the original workshop report and the Au, Robertson, and Shallit reference above. This is supplied dated curation derived from the OWR record, not an independently verified literature assessment; time-sensitive status and cited claims require independent verification.
\end{itemize}
\end{document}
